\documentclass[12pt,a4paper]{article}
\usepackage{ctex}
\usepackage{amsmath,amscd,amsbsy,amssymb,latexsym,url,bm,amsthm, amsfonts}
\usepackage{epsfig,graphicx,subfigure}
\usepackage{enumitem,balance}
\usepackage{wrapfig}
\usepackage{mathrsfs,euscript}
\usepackage[usenames]{xcolor}
\usepackage{hyperref}
\usepackage{tikz}
\usepackage[vlined,ruled,linesnumbered]{algorithm2e}
\hypersetup{colorlinks=true,linkcolor=black}
\ctexset{figurename=Figure}
\ctexset{tablename=Table}

\newtheorem{theorem}{Theorem}
\newtheorem{lemma}[theorem]{Lemma}
\newtheorem{proposition}[theorem]{Proposition}
\newtheorem{corollary}[theorem]{Corollary}
\newtheorem{exercise}{Exercise}
\newtheorem*{solution}{Solution}
\newtheorem{definition}{Definition}
\theoremstyle{definition}

\newcommand{\postscript}[2]
{\setlength{\epsfxsize}{#2\hsize}
\centerline{\epsfbox{#1}}}

\renewcommand{\baselinestretch}{1.0}

\setlength{\oddsidemargin}{-0.365in}
\setlength{\evensidemargin}{-0.365in}
\setlength{\topmargin}{-0.3in}
\setlength{\headheight}{0in}
\setlength{\headsep}{0in}
\setlength{\textheight}{10.1in}
\setlength{\textwidth}{7in}
\makeatletter \renewenvironment{proof}[1][Proof]
{\par\pushQED{\qed}\normalfont\topsep6\p@\@plus6\p@\relax\trivlist
\item[\hskip\labelsep\bfseries#1\@addpunct{.}]\ignorespaces}{\popQED\endtrivlist\@endpefalse}
\makeatother
\makeatletter
\renewenvironment{solution}[1][Solution]
{\par\pushQED{\qed}\normalfont\topsep6\p@\@plus6\p@\relax\trivlist
\item[\hskip\labelsep\bfseries#1\@addpunct{.}]\ignorespaces}{\popQED\endtrivlist\@endpefalse}
\makeatother

\begin{document}
\noindent

%========================================================================
\noindent\framebox[\linewidth]{\shortstack[c]{
    \Large{\textbf{Lab 02-Greedy Algorithm}}\vspace{1mm}\\
CS2312-Algorithm and complexity, Xiaofeng Gao, Spring 2026.}}
\begin{center}
  \footnotesize{\color{red}$*$ Contact TA Yixuan Cai for any
    questions. Include both your .pdf and .tex files in the uploaded
  .rar or .zip file.}

  % Please write down your name, student id and email.
  \footnotesize{\color{blue}$*$ Name:Hengtao Wu  \quad
  Student ID:524031910038 \quad Email: Eternity\_w@sjtu.edu.cn}

\end{center}

\begin{enumerate}
  \item {\color{purple}(Fractional Knapsack)} Let \((w_{1},v_{1}), (w_{2},v_{2}), \ldots, (w_{n},v_{n})\) be \(n\) pairs of positive real values. 
\begin{enumerate}
    \item Given a real value \(W \leq \sum_{i=1}^{n} w_{i}\), design an algorithm to find \(x_{1}, x_{2}, \ldots, x_{n}\) to maximize the objective function
\[
\sum_{i=1}^{n} \frac{v_{i}}{w_{i}} \cdot x_{i}
\]
subject to
\begin{itemize}
    \item $0 \leq x_{i} \leq w_{i} \quad \text{for every } i \in [1,n]$;
    \item $\sum_{i=1}^{n} x_{i} \leq W$.
\end{itemize}
    \item Analyze the time complexity of your algorithm.
\end{enumerate}

\begin{solution}
    \begin{enumerate}
        \item 
             \textbf{Algorithm Principle:} The algorithm greedily selects items with the highest ratio ($v_i/w_i$), taking them fully if capacity allows, or fractionally to fill the knapsack. \\
             The pseudo-code is presented below.

        \begin{minipage}{\linewidth}
            \begin{algorithm}[H]
            \setcounter{AlgoLine}{0}
            \SetKwInOut{Input}{Input}\SetKwInOut{Output}{Output}
            \Input{\((w_{1},v_{1}), (w_{2},v_{2}), \ldots, (w_{n},v_{n})\), $W$}
            \Output{$A = \{x_{1}, x_{2}, \ldots, x_{n}\}$}
            \BlankLine
            Initialize $x_i = 0$ for every $i \in [1,n]$\;
            Sort the items such that \(\frac{v_1}{w_1} \geq \frac{v_2}{w_2} \geq \ldots \geq \frac{v_n}{w_n}\)\;
            \For{\(i=1,2, \ldots, n\)}{
                \If{\(W==0\)}{
                    \textbf{break}\;
                }
                \eIf{\(W\geq w_i\)}{
                    \(x_i \leftarrow w_i\)\ \tcp*[r]{Greedy}
                    \(W \leftarrow W - w_i\)\;
                }{
                    \(x_i \leftarrow W\)\;
                    \(W \leftarrow 0\)\;
                    \textbf{break}\;
                }
            }
            \Return $A = \{x_{1}, x_{2}, \ldots, x_{n}\}$\;
            \caption{Fractional Knapsack Greedy Algorithm}
           \end{algorithm}

        \end{minipage}
        
            \textbf{Proof of Optimality (Exchange Argument):} \\
            Let the items be sorted by their value-to-weight ratio such that $\frac{v_1}{w_1} \geq \frac{v_2}{w_2} \geq \ldots \geq \frac{v_n}{w_n}$. Let $G = \{x_1, x_2, \ldots, x_n\}$ be the solution produced by our greedy algorithm, and $O = \{y_1, y_2, \ldots, y_n\}$ be an optimal solution. Assume, for the sake of contradiction, that $G \neq O$. 
            
            Let $i$ be the first index where the two solutions differ (i.e., $x_k = y_k$ for all $k < i$, and $x_i \neq y_i$). Because the greedy algorithm always takes as much of the highest-ratio available item as possible, we must have $x_i > y_i$. To compensate for taking less of item $i$, the optimal solution $O$ must have taken more of some other item $j$ (where $j > i$) compared to the greedy solution. Since the items are sorted, we know $\frac{v_i}{w_i} \geq \frac{v_j}{w_j}$.
            
            Consider a small weight $\Delta = \min(x_i - y_i, y_j) > 0$. If we remove $\Delta$ weight of item $j$ from $O$ and replace it with $\Delta$ weight of item $i$, the total weight of the knapsack remains valid and unchanged. However, the total value of this new solution changes by:
            $$\Delta \cdot \left( \frac{v_i}{w_i} - \frac{v_j}{w_j} \right) \geq 0$$
            
            This exchange yields a new valid solution that is at least as valuable as $O$. By repeatedly applying this exchange process for all differences, we can eventually transform the optimal solution $O$ entirely into the greedy solution $G$ without ever decreasing the total value. Therefore, the greedy solution $G$ must also be a optimum.
        
            
           
            
           \item Complexity Analysis

            \begin{itemize}
               \item \textbf{Sorting:} Calculating the fraction $\frac{v_i}{w_i}$ for each of the $n$ items takes $\mathcal{O}(n)$ time. Sorting these $n$ items in descending order based on their fractions takes $\mathcal{O}(n \log n)$ time.
               \item \textbf{Greedy Choice:} The algorithm iterates through the sorted items using a single \textbf{for} loop. In the worst case, it runs $n$ times. Inside the loop, all conditional checks take constant time $\mathcal{O}(1)$. Thus, this greedy selection phase takes $\mathcal{O}(n)$ time.
            \end{itemize}

            Overall,the time complexity is dominated by the sorting step, which yields a total time complexity of $\mathcal{O}(n \log n) + \mathcal{O}(n) = \mathcal{O}(n \log n)$.
    \end{enumerate}
\end{solution}




  \item {\color{purple} (Min Interval Covering)} Given a set of $n$ distinct points $\{x_1, x_2, \dots, x_n\}$ distributed on real number axis with $x_1 < x_2 < \dots < x_n$, cover these points using closed intervals of length $1$.  
\begin{enumerate}
    \item Design an algorithm to find the minimum number of such intervals and their positions.
    \item Provide a correctness proof of your algorithm.
\end{enumerate}

    \begin{solution}
        \begin{enumerate}
            \item Algorithm design.
            The algorithm iterates through the points from left to right. Whenever it encounters the first uncovered point $x_i$, it places an interval starting exactly at $x_i$, extending to $x_i + 1$. This greedily covers as many subsequent points as possible before needing a new interval.

            \begin{minipage}{\linewidth}
                \begin{algorithm}[H]
                \setcounter{AlgoLine}{0}
                \SetKwInOut{Input}{Input}\SetKwInOut{Output}{Output}
                \Input{A set of distinct points $\{x_1, x_2, \ldots, x_n\}$ with $x_1 < x_2 < \ldots < x_n$}
                \Output{Minimum number of intervals $count$, and their positions $I$}
                \BlankLine
                Initialize $I = \emptyset$\,$count = 0$\,$i = 1$\;
                \While{$i \leq n$}{
                    Add the interval $[x_i, x_i + 1]$ to $I$\;
                    \(count \leftarrow count + 1\)\;
                    \tcp{Find the first point not covered by $[x_i, x_i + 1]$}
                    Let $x_k$ be the smallest point such that $x_k > x_i + 1$\;
                    \eIf{such $x_k$ exists}{
                        $i \leftarrow k$\;
                    }{
                        \textbf{break}\;
                    }
                }
                \Return $count$, and $I$\;
                \caption{Greedy Min Interval Covering}
                \end{algorithm}
            \end{minipage}

            \item
            \textbf{Theorem.}Greedy Min Interval Covering algorithm is optimal.\\
            \textbf{Proof.}
                Let $G = \{g_1, g_2, \ldots, g_k\}$ be the sequence of intervals chosen by our greedy algorithm, and $O = \{o_1, o_2, \ldots, o_m\}$ be the intervals of an optimal solution.
        
                Assume, for the sake of contradiction, that $G \neq O$. Let $j$ be the first index where the two solutions differ (i.e., $g_r = o_r$ for all $r < j$, but $g_j \neq o_j$). 
                
                Let $x^*$ be the first point not covered by the first $j-1$ intervals. By our greedy strategy, the algorithm places $g_j = [x^*, x^* + 1]$.Since the optimal solution $O$ is valid, it must also cover $x^*$. Therefore, its $j$-th interval must be $o_j = [y, y + 1]$ where $y \leq x^*$. Because $y \leq x^*$, it strictly implies that the right endpoint $y + 1 \leq x^* + 1$. 
                
                This means the right endpoint of the greedy interval $g_j$ extends at least as far to the right as the optimal interval $o_j$. Hence, $g_j$ covers all the unhandled points that $o_j$ covers. 
                
                If we exchange $o_j$ with $g_j$ in the optimal solution $O$, no points will be left uncovered.This exchange creates a new valid solution that uses the exact same number of intervals as $O$ (thus still optimal), but matches $G$ up to the $j$-th interval.
                
                By inductively repeating this exchange process for all differences, we can transform the optimal solution $O$ entirely into the greedy solution $G$ without ever increasing the number of intervals. Therefore, the greedy algorithm produces a minimum set of intervals.
            
        \end{enumerate}
    \end{solution}


  \item {\color{purple} (Matroid)} Consider a pair $M = (U, \mathcal{I})$ where universe $U$ is a finite set and $\mathcal{I}
    \subseteq 2^U$ is a collection of subsets of $U$. We call $M$ a
    matroid if it satisfies \emph{hereditary property} and
    \emph{exchange property}:
    \begin{itemize}
      \item (\emph{HEREDITARY PROPERTY}) $\mathcal{I}$ is nonempty
        and for every $A \in \mathcal{I}$ and $B \subseteq A$, it
        holds that $B \in \mathcal{I}$.
      \item (\emph{EXCHANGE PROPERTY}) For any $A, B \in
        \mathcal{I}$ with $|A| < |B|$, there
        exists some $x \in B \setminus A$ such that $A \cup \{x\} \in
        \mathcal{I}$.
    \end{itemize}
    Each set $A \in \mathcal{I}$ is called an \emph{independent set}. The maximal independent sets are called \emph{bases}.
    \begin{enumerate}
      \item The dual of a matroid $\mathcal{I}$ over a universe $U$ is the family of subsets $\mathcal{I}'$ over the same universe $U$, consisting of the complements of bases of $\mathcal{I}$, and all subsets of these sets. Prove that the dual $\mathcal{I}'$ is itself a matroid.
      \item Consider a matroid $\mathcal{I}$ over a universe $U$, and for a parameter $k$, define the following family of subsets 
      \[
          \mathcal{I}_k = \{A \mid A \in I \,\, \text{and} \,\, |A| \leq k\}
      \]
      Prove that for each positive integer $k$, $\mathcal{I}_k$ is a matroid.
    \end{enumerate}

    \begin{solution}
        \begin{enumerate}
    
            \item 
            \begin{proof}
We demonstrate that the dual $\mathcal{I}'$ is a matroid by verifying the hereditary property and the exchange property.

    \begin{itemize}
        \item (\emph{HEREDITARY PROPERTY})
        Let $A \in \mathcal{I}'$ and $A' \subseteq A$. By the definition of the dual matroid, there exists a base $B$ of $\mathcal{I}$ such that $A$ is a subset of the complement of $B$, i.e., $A \subseteq U \setminus B$. Since $A' \subseteq A$, it follows directly that $A' \subseteq U \setminus B$. Thus, $A'$ is also a subset of the complement of a base, meaning $A' \in \mathcal{I}'$.


    \item (\emph{EXCHANGE PROPERTY})
    Let $A, B \in \mathcal{I}'$ with $|A| < |B|$. We need to find an element 
    $x \in B \setminus A$ such that $A \cup \{x\} \in \mathcal{I}'$.

    Since $A, B \in \mathcal{I}'$, there exist bases $B_A, B_B$ of $\mathcal{I}$ 
    such that $B_A \cap A = \emptyset$ and $B_B \cap B = \emptyset$.

    To prove $A \cup \{x\} \in \mathcal{I}'$, it suffices to construct a base 
    $\hat{B}$ of $\mathcal{I}$ satisfying $\hat{B} \cap A = \emptyset$ and 
    $\hat{B} \cap \{x\} = \emptyset$ for some $x \in B \setminus A$.

    Let $B_B \cap A = \{e_1, e_2, \ldots, e_k\}$. We construct $\hat{B}$ from 
    $B_B$ by iteratively replacing each $e_i \in A$ with a safe element from 
    $B_A$. Specifically, by the base exchange property of matroids, for each 
    $e_i$, there exists $f_i \in B_A$ such that $(B_B \setminus \{e_i\}) \cup 
    \{f_i\}$ is again a base. Applying this exchange $k$ times yields a base 
    $\hat{B}$ with $\hat{B} \cap A = \emptyset$, since all elements of 
    $B_B \cap A$ have been removed and replaced by elements of $B_A \subseteq 
    U \setminus A$.

    It remains to show that $\hat{B}$ does not contain all elements of 
    $B \setminus A$. Since $B_B \cap B = \emptyset$, the only elements of $B$ 
    that could appear in $\hat{B}$ are those introduced via the $f_i$'s. 
    Therefore:
    $$|\hat{B} \cap (B \setminus A)| \leq k = |B_B \cap A|$$

    Since $B_B \cap B = \emptyset$, every element of $B_B \cap A$ belongs to 
    $A \setminus B$, so $k \leq |A \setminus B|$. Combined with the assumption 
    $|A| < |B|$, which implies $|A \setminus B| < |B \setminus A|$, we obtain:
    $$|\hat{B} \cap (B \setminus A)| \leq |A \setminus B| < |B \setminus A|$$

    Hence $\hat{B}$ does not contain every element of $B \setminus A$. There 
    exists at least one $x \in B \setminus A$ with $x \notin \hat{B}$. Since 
    $\hat{B} \cap A = \emptyset$ and $x \notin \hat{B}$, we have 
    $\hat{B} \subseteq U \setminus (A \cup \{x\})$, which gives 
    $A \cup \{x\} \in \mathcal{I}'$.
        
    \end{itemize}
\end{proof}
    
            \item 
            \begin{proof}
                We verify the two properties for $\mathcal{I}_k = \{A \mid A \in \mathcal{I} \text{ and } |A| \leq k\}$.
                
                \begin{itemize}
                    \item (\emph{HEREDITARY PROPERTY})
                    Let $A \in \mathcal{I}_k$ and $B \subseteq A$. Since $A \in \mathcal{I}_k$, we know $A \in \mathcal{I}$ and $|A| \leq k$. Because $\mathcal{I}$ is a matroid, the hereditary property guarantees that any subset of an independent set is independent, so $B \in \mathcal{I}$. Furthermore, since $B \subseteq A$, we have $|B| \leq |A| \leq k$. Thus, $B \in \mathcal{I}$ and $|B| \leq k$, meaning $B \in \mathcal{I}_k$.
                    
                    \item (\emph{EXCHANGE PROPERTY})
                    Let $A, B \in \mathcal{I}_k$ with $|A| < |B|$. Since $\mathcal{I}_k \subseteq \mathcal{I}$, both $A$ and $B$ are independent sets in the original matroid $\mathcal{I}$. By the exchange property of $\mathcal{I}$, there exists an element $x \in B \setminus A$ such that $A \cup \{x\} \in \mathcal{I}$.
                    
                    We only need to verify the size constraint for $A \cup \{x\}$. Since $|A| < |B|$ and both are integers, we have $|A| \leq |B| - 1$. We also know that $B \in \mathcal{I}_k$, which implies $|B| \leq k$. 
                    Substituting this gives $|A| \leq k - 1$. 
                    Therefore, the size of the new set is:
                    $$ |A \cup \{x\}| = |A| + 1 \leq (k - 1) + 1 = k $$
                    Since $A \cup \{x\} \in \mathcal{I}$ and $|A \cup \{x\}| \leq k$, we conclude that $A \cup \{x\} \in \mathcal{I}_k$.
                \end{itemize}
            \end{proof}
        
        \end{enumerate}
    \end{solution}
% \item{\color{purple} Discontinuous Classroom Assignment} Recall the classroom scheduling example on class. This time, however, a lecture is no longer continuous. Instead, it is broken into several timepieces, namingly $\{(ts_i^1,te_i^1),(ts_i^2,te_i^2),...\}$ for lecture $i$, where $ts_i^j$ and $te_i^j$ stands for the start and ending time for timepiece $j$. You still need to arrange all the timepieces of a single lecture to the same classroom, and avoid conflict.
% \begin{enumerate}
%     \item Is the number of classroom needed still equals to the depth? Prove your answer. 
%     \item Design an algorithm to find the minimal number of classroom needed.
% \end{enumerate}
\end{enumerate}

\newpage
        \noindent\rule{\textwidth}{0.4pt} 
        \vspace{0.2cm}
        
        \textbf{\large Appendix: An Alternative Approach (For TA's Feedback)} \\
        \textit{Dear TA, the proof above is my final answer. I actually tried another method first (the "capacity" argument below), but I wasn't 100\% sure about it. Could you help check if this alternative proof is also acceptable, or if I missed something subtle? Really appreciate your feedback!}

        \begin{proof}[Alternative Proof]
        We demonstrate the exchange property by building a new base and analyzing its capacity.
        Let $A, B \in \mathcal{I}'$ with $|A| < |B|$. We need to find an element $x \in B \setminus A$ such that $A \cup \{x\} \in \mathcal{I}'$.

        Since $A, B \in \mathcal{I}'$, there exist bases $B_A, B_B$ of $\mathcal{I}$ such that $A \subseteq U \setminus B_A$ and $B \subseteq U \setminus B_B$. Taking the complement of both sides gives $B_A \subseteq U \setminus A$ and $B_B \subseteq U \setminus B$.

        Consider the set $I = B_B \setminus A$. Since $I \subseteq B_B \in \mathcal{I}$, $I$ is an independent set in $\mathcal{I}$. Moreover, $I \subseteq U \setminus A$ by construction.

        We claim we can extend $I$ to a base $B^*$ of $\mathcal{I}$ that is entirely contained in $U \setminus A$. To see this, note that $B_A \subseteq U \setminus A$ and $B_A$ is a base of $\mathcal{I}$, so $r(U \setminus A) = r(U) = r$ (the rank of the whole matroid). Since $I \subseteq U \setminus A$ is independent, we can extend $I$ to a maximal independent set $B^*$ within $U \setminus A$. Since $r(U \setminus A) = r$, this maximal independent set is a base of $\mathcal{I}$, and by construction $B^* \subseteq U \setminus A$, so $B^* \cap A = \emptyset$.

        Now, assume for the sake of contradiction that for all $x \in B \setminus A$, $A \cup \{x\} \notin \mathcal{I}'$. This implies that no element in $B \setminus A$ lies outside of $B^*$, i.e., $B \setminus A \subseteq B^*$.

        Since $I = B_B \setminus A \subseteq U \setminus B$, $I$ contains no elements from $B$. Since $B \setminus A \subseteq B^*$ and $B^* \subseteq U \setminus A$, all elements of $B \setminus A$ lie in $B^* \setminus I$, so:
        $$|B \setminus A| \leq |B^*| - |I| = r - |B_B \setminus A|$$

        Since $B_B \subseteq U \setminus B$, we have $B_B \cap B = \emptyset$, so $B_B \cap A = B_B \cap (A \setminus B)$, which gives $|B_B \cap A| \leq |A \setminus B|$.

        Therefore:
        $$|B \setminus A| \leq r - |B_B \setminus A| = r - (|B_B| - |B_B \cap A|) = |B_B \cap A| \leq |A \setminus B|$$

        Using the fact that $|X \setminus Y| = |X| - |X \cap Y|$, we get:
        $$|B| - |A \cap B| \leq |A| - |A \cap B| \implies |B| \leq |A|$$
        This directly contradicts our initial assumption that $|A| < |B|$.

        Therefore, our assumption must be false. There must exist at least one $x \in B \setminus A$ such that $x \notin B^*$. Since $x \in B \setminus A$, we have $x \notin A$. Combined with $x \notin B^*$, we get $B^* \cap (A \cup \{x\}) = \emptyset$, i.e., $B^* \subseteq U \setminus (A \cup \{x\})$, which implies $A \cup \{x\} \in \mathcal{I}'$.
        
        \end{proof}


%========================================================================
\end{document}

